\section{Introduction}
Rapid urbanization is altering the timing of plant life-cycle events—phenology—in cities, with important implications for urban forestry and localized ecosystem services, including canopy cooling, pollen exposure, carbon uptake, and wildlife interactions \citep{Yang2023GCB, Zhou2022NCC}. Evidence for urban phenological shifts includes earlier spring leaf-out and later fall senescence, patterns often attributed to elevated urban temperatures associated with the urban heat island effect \citep{Kabano2021LUP, Zhang2006JGRB}. Most existing evidence, however, stems from urban-rural comparisons or citywide averages, which implicitly treat each city as phenologically uniform \citep{Kabano2021LUP, Li2017GCB, Meng2020PNAS, Wang2019NatEcoEvol, Zhang2006JGRB}. This assumption is unlikely to hold in practice. Cities are spatially heterogeneous mosaics of built surfaces, vegetation, management regimes, and planted species. As a result, individual trees within the same city may experience different local thermal environments and exhibit different phenological timing \citep{Wu2025NatCommun}. Ignoring this intra-city heterogeneity can obscure the temperature-phenology relationships that operate within cities and limit the use of cities as natural laboratories for understanding ecological responses to climate change \citep{Zhou2022NCC}. It can also hinder urban-forest management by masking neighborhood-scale variation in ecosystem services and disservices that affect public health and well-being \citep{Crawford2024ERL, Katz2019STOTEN, Song2025SRS}. 

Increasing evidence shows that phenological timing can vary markedly among individual trees as a function of the local environments they experience \citep{Christmann2023UrbanEcosyst, Wohlfahrt2019NatEcoEvol, Wu2024NCC, Zellweger2020Science}. In cities, two sources of variation are likely to be especially important: fine-scale climate variation and species diversity. First, individual trees experience substantial intra-city variation in air temperature, a major environmental cue for plant phenology \citep{penuelas_phenology_2009, piao_plant_2019, Wohlfahrt2019NatEcoEvol}. This thermal heterogeneity arises from spatial variation in urban form, surface materials, vegetation cover, and anthropogenic heat emissions \citep{Stewart2012BAMS}. Reported intra-city seasonal temperature differences can reach up to 5-6 $^\circ$C, comparable to or greater than many urban-rural contrasts \citep{Cao2021LUP, Hart2009TAC, Johnson2020UrbanClimate, Kong2025ERL, Kousis2021ScientificReports, Naserikia2023STOTEN, Venter2020RSE}. Second, urban trees are composed of diverse assemblages of intentionally planted species with distinct phenological schedules and environmental sensitivities \citep{Gasparovic2023Forests, Parece2018UrbanSci}. Because multiple species are often planted across the same neighborhoods and streets, species composition can amplify observed phenological heterogeneity and motivate analyses that explicitly link individual-level phenology to taxonomic identity.

Although single-city studies have documented intra-urban variation in tree phenology, whether these patterns generalize across broader climatic contexts remains unclear \citep{Crawford2024ERL, Katz2019STOTEN, Parece2018UrbanSci, Zipper2016ERL, Wu2025NatCommun}. Regional climate may modify temperature-phenology relationships through chilling limitation, heat stress, moisture constraints, and differences in species composition or provenance \citep{Chen2020NCC, Chuine2010Science, Jochner2015EnvironmentalPollution, Fu2015Nature, Li2019NatEcoEvol, Zhou2022NCC}. However, it remains poorly understood how these broader climatic constraints interact with fine-scale intra-urban thermal heterogeneity. The contiguous United States provides a useful setting for addressing this issue because it spans large gradients in climate, urban form, and tree species composition. At the same time, the availability of municipal tree inventories, high-resolution satellite imagery, and gridded urban air-temperature products enables analysis of phenology at the level of individual trees across many cities. 

Here, we examine individual-tree phenology across 76 cities in the contiguous United States by integrating municipal street-tree inventories, PlanetScope satellite time series, and 1-km near-surface urban air-temperature data. We address three questions. First, how large is intra-city variation in tree phenology relative to variation among cities? We retrieve start-of-season (SOS) and end-of-season (EOS) metrics for more than 1.6 million inventoried deciduous street trees from 345 species and partition phenological variation into intra-city, cross-city, and interannual components. We expect substantial intra-city variation because individual trees experience different local environments and urban tree communities have diverse species. Second, how is individual-tree phenology associated with fine-scale intra-city temperature variation? For each species within each city, we estimate the association between phenological timing and local preseason air temperature while accounting for precipitation, radiation, spatial autocorrelation, and interannual variability. We expect warmer local environments to be associated with earlier SOS, later EOS, and consequently longer growing-season-length (GSL) for many species, while also expecting substantial variation in magnitude and direction among species–city combinations. Third, do these temperature-phenology associations vary with regional climatic context? We test whether estimated temperature sensitivities vary with city-level mean annual temperature and precipitation. We expect apparent temperature sensitivities to weaken or reverse in hotter or drier cities, where chilling limitation, heat stress, or moisture constraints may limit phenological advancement or growing-season extension.